\section{Supplement to Section 4: finite-amplitude instability}\label{supp:bistab}

This section gives the protocol details and secondary results of Section~\ref{M-sec:bistability}.

\subsection{Entries, kicks and surface memory}\label{supp:bistab-prot}

In the face entry a point of the cutting edge at angle $\phi$ is in the material once the tool centre has advanced by $R(1-\sin\phi)$, so that full engagement is reached after $R/f_z=160$ teeth (40 revolutions). A depth ramp was not used, because a real axial ramp-in shifts the regeneration delay through the helix. The depth ladder of Section~\ref{M-sec:bistab-prot} starts at $0.25\alim$ and rises in steps of \SI{8}{\percent} up to $1.05\alim$; the first interval that ends in chatter is bisected to a relative \SI{0.3}{\percent}.
% src: src/tdsim.c (opt[1]-opt[3]: ramp_f, f0, edge), src/tdsim.py (entry_revs), src/c1_unsafe_zone.py (ENTRIES, LADDER: growth 1.08, lo 0.25, hi 1.05), src/tdboundary.py (boundary: rtol 3e-3); README §11.2

The critical kick was searched upwards on a grid of 16 kicks, $\SI{0.1}{\milli\metre\per\second}\times2^k$ ($k=0,\dots,14$) and \SI{3}{\metre\per\second} (a swing of about \SI{3.7}{\milli\metre}), and only the first interval that ended in chatter was bisected, six times in $\log\Delta v$ (resolution \SI{1.1}{\percent}). The lower edge of the unsafe band was located by six bisection steps on $[0.05,0.98]\,\alim$; at each step seven kicks, spaced geometrically from \SI{3}{\milli\metre\per\second} to \SI{3}{\metre\per\second}, were applied at two phases half a tooth period apart. The bisection assumes that the band has no holes in depth. At $0.98\alim$ the slowest mode decays by only 0.985--0.994 per tooth period, so a residual of the face entry at full engagement (revolution 40) decays to 3--\SI{22}{\percent} of its value by the kick at revolution 100. Measured on the unkicked run, the deviation from the forced orbit is at most \SI{5.2}{\nano\metre} at revolution 40, \SI{0.85}{\nano\metre} at revolution 100 and \SI{0.01}{\nano\metre} at revolution 300 (Section~\ref{M-sec:bistab-checks}).
% src: src/c1_unsafe_zone.py (DV_GRID = 1e-4 x 2^k, k = 0..14, and 3 m/s; N_BISECT = 6; BAND_DV = geomspace(3e-3, 3, 7); BAND_PHASES = 2 at offsets 0 and ND/2; BAND_BISECT = 6 on [0.05, 0.98]; KICK_REV = 100); 3 m/s / omega_n = 3.69 mm; 2^(1/64) = 1.011; results/c1_basin.npz (decay at 0.98 a_lim: 0.9937, 0.9854, 0.9933, 0.9925 at 1090, 1420, 2500, 2870 rpm; decay^240 = 0.22, 0.029, 0.20, 0.16 over the 60 revolutions = 240 tooth periods between revolutions 40 and 100)

A kick displaces the workpiece by about $\Delta v/\omega_n$, \SI{37}{\micro\metre} for \SI{30}{\milli\metre\per\second}, less than one feed, but the clipped chatter that it releases reaches 0.06--\SI{2}{\milli\metre} peak to peak (Section~\ref{M-sec:bistab-attr}), up to some forty feeds, and a kick of \SI{3}{\metre\per\second} alone swings the workpiece by about 74 feeds, so the tool leaves the cut for many passes. Four remembered passes cover only \SI{0.2}{\milli\metre} of feed; with them, two of the four clipped-chatter cases of \ref{M-app:simulator} exceed the \SI{1}{\centi\metre} at which the simulator stops where the full memory stays bounded, and a third is \SI{27}{\percent} too large. Even with the full memory a run can exceed \SI{1}{\centi\metre}, because the chip is not limited when the tool plunges into the material, for instance the step entry at \SI{1050}{\rpm} and $0.9\alim$, where the brute-force reference does so as well; with the stop raised to \SI{10}{\centi\metre}, that run returns to the forced orbit after an excursion of \SI{15}{\milli\metre} (Section~\ref{supp:bistab-attr}).
% src: verification/tdsim/check9_full_memory.json (bounded chatter 0.4345 and 0.5896 mm; nreg 4: 2350 rpm and 1090 rpm diverge where full memory is bounded, 2500 rpm 0.554 vs 0.435 mm; 1050 rpm, a = 24.18 mm = 0.9 a_lim, diverges with full memory and with the 600-pass brute force); 30 mm/s / omega_n = 36.9 um, 3 m/s / omega_n = 3.69 mm = 74 f_z; 4 f_z = 0.2 mm; README §11.2

\subsection{Two outcomes and entry limits}\label{supp:bistab-map}

In Fig.~\ref{M-fig:two} the command of the forced orbit lies between $-6.24$ and \SI{5.49}{\volt}; after the kick of \SI{7.6}{\milli\metre\per\second} the displacement settles between $-393$ and \SI{356}{\micro\metre}, with the command switching between $\pm\Vmax$. Over the 16 phases the critical kick at this depth is \SIrange{7.2}{8.8}{\milli\metre\per\second}. After the kick of \SI{7.0}{\milli\metre\per\second} the command still reaches the clip repeatedly while the deviation decays. A kick of \SI{30}{\milli\metre\per\second} on the settled cut at \SI{15}{\milli\metre} ends in chatter with the clip but decays without it.
% src: paper/figures/figure_numbers.json (fig03: orbit_V -6.24..5.49 V, chatter_x_um -393..356, chatter_V -9..9 V); results/c1_basin.npz (2500 rpm, 0.9 a_lim, clipped: range 7.21-8.76 mm/s); README §11.3

Where the step entry binds, the better of the face entry and the feed ramp gains \SI{2.51}{\milli\metre} on average; the largest gain, \SI{4.74}{\milli\metre}, is at \SI{1980}{\rpm} (\SI{19.0}{\milli\metre} for the step entry against \SI{23.8}{\milli\metre}), and at \SI{2500}{\rpm} the limits are 13.0 and \SI{17.25}{\milli\metre}. At 500, 540, 670, 720, 1000 and \SI{1080}{\rpm} the face entry and the feed ramp, and at 500, 550, 670 and \SI{1000}{\rpm} the clipped step entry, reached the top of the depth ladder, $1.05\alim$; at the other 245 speeds the mean face-entry limit is \SI{18.20}{\milli\metre} against a mean $\alim$ of \SI{18.21}{\milli\metre}. Without the clip the step-entry limit exceeds $\alim$ by more than \SI{1.5}{\percent} at 75 speeds, among them every speed from 1990 to \SI{2480}{\rpm}, and by at least \SI{5}{\percent}, the top of the depth ladder, at \SI{2000}{\rpm}. Near the forced resonance (\SIrange{1970}{2030}{\rpm}) runs of 2000 revolutions place the onset at 1.004--$1.012\alim$, so the excess there is the bias of the classifier; at the other speeds it was not checked with long runs. No unsafe band was found at 17 speeds: 500, 540 and \SI{720}{\rpm} ($g=\SI{1500}{\volt\second\per\metre}$) and 1660--1800 and \SIrange{1840}{1940}{\rpm}, where the gain is \SIrange{135}{979}{\volt\second\per\metre}.
% src: results/c1_summary.json (map: soft_minus_step_mean_mm 2.51, at_2500: step_clip 13.0, edge_clip 17.25); results/c1_map.npz recomputed: largest gain 4.741 mm at 1980 rpm (a_step_clip 19.02 mm, max(a_edge_clip, a_feed_clip) 23.76 mm); a_edge_clip and a_feed_clip = 1.05 a_lim at 500, 540, 670, 720, 1000, 1080 rpm, a_step_clip at 500, 550, 670, 1000 rpm; a_step_free/a_lim > 1.015 at 75 speeds (1.0165-1.050, all of 1990-2480 rpm, 1.050 at 2000 rpm); verification/tdsim/check10_forced_resonance_long_runs.json (onset 1.0039-1.0121 at 1970-2030 rpm); results/c1_summary.json (band: none_at_rpm, g_where_none 135-979 V s/m)

\subsection{Basin without the clip and near the limit}\label{supp:bistab-basin}

Worst-phase kicks of at most \SI{10}{\milli\metre\per\second} end in chatter from 0.95, 0.90 and $0.95\alim$ upwards at 1420, 2500 and \SI{2870}{\rpm} and at none of the 12 depths at \SI{1090}{\rpm}; kicks of at most \SI{30}{\milli\metre\per\second} from 0.40, 0.80, 0.40 and $0.50\alim$ upwards at 1090, 1420, 2500 and \SI{2870}{\rpm}. Over the 16 phases the critical kick varies by at most \SI{15}{\percent} at $0.98\alim$ and \SI{24}{\percent} at any depth.
% src: results/c1_basin.npz recomputed: smallest depth with worst-phase dv <= 10 mm/s: 0.95, 0.90, 0.95 a_lim at 1420, 2500, 2870 rpm, none at 1090; <= 30 mm/s: 0.40, 0.80, 0.40, 0.50 a_lim at 1090, 1420, 2500, 2870 rpm; results/c1_summary.json (basin: phase_spread_098, phase_spread_max)

Without the clip, kicks of 0.78, 3 and \SI{1.64}{\metre\per\second}, each in one phase of 16, ended in chatter at \SI{1090}{\rpm} and 0.9 and $0.98\alim$ and at \SI{2500}{\rpm} and $0.98\alim$; there the clip shrinks the basin by factors of 52, 276 and 701. At nine further points the clipped cut chatters but the unclipped cut did not for any kick up to \SI{3}{\metre\per\second}, so the factor is at least 70--595 (Fig.~\ref{M-fig:basin}). At the four remaining, shallow points (\SI{1420}{\rpm} at 0.3 and $0.6\alim$; 2500 and \SI{2870}{\rpm} at $0.3\alim$) no kick up to \SI{3}{\metre\per\second} ends in chatter, with or without the clip.
% src: results/c1_basin.npz (vmax = 0 rows: finite minima 0.784, 3.000, 1.638 m/s, one phase each; ratios of the smallest unclipped to the smallest clipped kick 52.1, 275.7, 700.9; lower bounds 3 m/s / clipped minimum 69.7-594.9 at the nine points listed); README §11.5

At $0.999\alim$, with the same protocol at four phases a quarter of a tooth period apart, the critical kick is \SIrange{1.41}{1.47}{\milli\metre\per\second} at \SI{2500}{\rpm}, \SIrange{7.1}{7.6}{\milli\metre\per\second} at \SI{2870}{\rpm} and \SIrange{9.1}{10.5}{\milli\metre\per\second} at \SI{1090}{\rpm}; \SI{1420}{\rpm} was not simulated at this depth.
% src: verification/tdsim/check11_kick_near_limit.json (phases 0, 4, 8, 12 of 16; 1.405-1.467, 7.055-7.611, 9.150-10.533 mm/s)

\subsection{Attractors beyond the basin}\label{supp:bistab-attr}

\paragraph{Long runs after a kick} At the 29 clipped points of Fig.~\ref{M-fig:basin} with a finite critical kick, the settled cut was kicked at revolution 100 of a face entry, at the phase with the smallest stored critical kick $\Delta v^*$, with 2, 2.5 and 3 times $\Delta v^*$ (at most \SI{3}{\metre\per\second}), and run for 2000 revolutions after the kick with the full surface memory. A run returned to the forced orbit if $\max\lvert x(t)-x(t-T)\rvert<0.01f_z$ over its last quarter; it ended in bounded, stationary chatter if, between revolutions 500--1000 and 1500--2000, the root mean square of $x$ and of $V$ differed by less than \SI{10}{\percent}, the fraction of clipped time by less than 0.05 and the dominant non-synchronous frequency by less than \SI{1}{\hertz}; it was stopped if $\lvert x\rvert$ exceeded \SI{1}{\centi\metre}. All 87 runs ended in bounded, stationary chatter, which the 300-revolution classifier had also called chatter. The four points at $0.9\alim$ were continued to 4000 revolutions.

\paragraph{Bifurcation-type diagrams} At 1090, 1420, 2500 and \SI{2870}{\rpm}, with the gain $g=\SI{1500}{\volt\second\per\metre}$, 28 depths evenly spaced from $\alim(0)$ to $1.05\alim$ were started by a face entry, by the face entry with a kick of 0.03 or \SI{0.3}{\metre\per\second} at revolution 100 (start of a tooth period), and by a step entry, and run for 1600 revolutions. The final state over the last 300 revolutions was the forced orbit, bounded chatter (stationary to \SI{10}{\percent}), a transient or an excursion beyond \SI{1}{\centi\metre}; only the first two occurred. The lowest grid depth reached by chatter was refined by five bisection steps against the depth below it, at each step with the seven kicks of the band study; as reaching the chatter branch from one of these starts is sufficient but not necessary, the refined end is an upper bound of the end of the branch.

\paragraph{Regularity} Twin runs whose kicks (four points at $0.9\alim$, 2000 revolutions) or depths (step entries at $1.05\alim$ at the four speeds and at two further cases, 1600 revolutions) differ by a relative $10^{-9}$ stay within $5\times10^{-7}$ of the peak-to-peak displacement of the chatter at the end, and their separation does not grow exponentially at any point: the chatter is regular, quasi-periodic where its Poincar\'e section is a closed curve.

\paragraph{Spot check} At the speeds of the band map with a band edge and $\alim\ge\SI{21}{\milli\metre}$, and at \SI{1050}{\rpm}, where the check of the surface memory had exceeded \SI{1}{\centi\metre} (700, 960, 980, 1000, 1020, 1040, 1050, 1060, 1080, 1620, 1640, 1820, 1960 and \SI{1980}{\rpm}, with the gains of Section~\ref{M-sec:feasibility}), the step entry and kicks of 0.03 and \SI{0.3}{\metre\per\second} at the start of revolution 100 of a face entry were run for 1600 revolutions at 0.9 and $0.98\alim$, and at $0.99\alim$ where the band edge lies above $0.9\alim$, at depths above the band edge only. Of 78 runs, 34 ended on the forced orbit, 23 in bounded chatter of 0.04--\SI{2.9}{\milli\metre} peak to peak, and 21 diverged, all of the latter at 700 and 1020--\SI{1060}{\rpm}.
% src: verification/paper/attractors_kick.py/.json (protocol, summary), attractors_bifurcation.py/.json/.npz (protocol, refine_lo/hi), attractors_chaos.py/.json (10 cases regular; sep_final_over_pp <= 4.9e-7), attractors_band_spot.py/.json (14 speeds, 78 runs: forced 34, chatter 23, diverged 21; bounded_pp_x_mm 0.044-2.94; speeds_with_divergence 700, 1020-1060)

\subsection{Checks of the protocols}\label{supp:bistab-checks}

\paragraph{Kick timing} At $0.98\,\alim$, where the chatter mode decays slowest, the face entry leaves at revolution 100 a deviation from the forced orbit of at most \SI{0.85}{\nano\metre} and \SI{0.73}{\micro\metre\per\second}, three orders of magnitude below the smallest critical kick. Kicking at revolution 300 instead of 100, at four phases a quarter of a tooth period apart at each of the four speeds, changes the critical kick at one of the 16 speed--phase pairs, by \SI{1.1}{\percent} (one bisection step), and its minimum over the phases at none.

\paragraph{Window} Judged over 900 instead of 300 revolutions after the kick, the 53 runs of the worst-phase searches at $0.98\,\alim$ get the same verdicts, and the critical kick is unchanged at the four speeds. The bias of the 300-revolution window found for step-entry limits near the forced resonance (Section~\ref{M-sec:bistab-prot}) does not reach the kick-based results, because after a kick the chatter grows or decays within some tens of revolutions (Fig.~\ref{M-fig:attr}g).

\paragraph{Phases} The critical kick at $0.98\,\alim$ was searched at the 48 phases between those of the stored campaign, i.e.\ at 64 phases per tooth period. The minimum over the phases falls by \SI{1.1}{\percent} at \SI{1090}{\rpm} (from 10.88 to \SI{10.76}{\milli\metre\per\second}, a value already reached with 32 phases) and is unchanged at the other three speeds; the spread over the phases grows from 13.9 to \SI{15.1}{\percent} at \SI{1090}{\rpm} and is unchanged elsewhere.

\paragraph{Holes in the band} At nine speeds spread over the range of the band (640, 960, 1640, 1820, 1000, 2140, 2800, 1300 and \SI{520}{\rpm}; first those where the stored edge jumps between neighbouring speeds or lies at the top of the bisection interval, then evenly spread ones), with the gains of the band map, each depth $0.05,0.10,\dots,0.95$ and $0.98\,\alim$ was tested exactly as one bisection step of the band protocol: the seven kicks from \SI{3}{\milli\metre\per\second} to \SI{3}{\metre\per\second} at two phases half a tooth period apart, stopping at the first chatter; 2028 runs in all. At every speed the lowest chattering grid depth lies directly above the bisection bracket of the stored edge, or inside it at 1000, 1640 and \SI{1820}{\rpm}, and no grid depth below the edge chatters. At eight speeds the chattering depths form one interval up to $0.98\,\alim$. At \SI{520}{\rpm} ($g=1500$, $\alim=\SI{18.4}{\milli\metre}$, edge $0.30\,\alim$) they do not: kicks end in chatter at 0.30--$0.55\,\alim$ and at 0.95 and $0.98\,\alim$, but at none of the seven grid depths from 0.60 to $0.90\,\alim$. The bisection of the band map assumes one interval; where this fails, the mapped edge is still the lowest depth at which a kick of the protocol ends in chatter, but the band above it is not unsafe throughout. Whether other kicks or phases would end in chatter inside the hole was not examined.
% src: verification/paper/protocol_band_holes.py/.json (protocol: SPEEDS order, FRACS 0.05:0.05:0.95 and 0.98, BAND_DV and BAND_PHASES of src/c1_unsafe_zone.py, face entry, kick at revolution 100, 300-revolution window, full memory; summary.*: grid_edge_frac 0.80/0.90/0.98/0.98/0.35/0.45/0.45/0.60/0.30 against stored_edge_frac 0.791/0.893/0.980/0.980/0.355/0.442/0.428/0.588/0.297; ambiguous_in_bracket at 1640, 1820, 1000 rpm; chatter_below_edge [] at all; chatter_set_is_interval_on_grid false only at 520 rpm: chatter_fracs 0.30-0.55, 0.95, 0.98, holes_above_edge 0.60-0.90; a_lim 18.45 mm; total_runs 2028)
% src: verification/paper/protocol_kick_timing.py/.json (units.residual_*: dev_x_from_orbit_um at rev100 <= 0.00085, dev_v_from_orbit_mmps <= 0.00073; summary.all: 16 points, ratio 0.989-1, max_abs_rel_change 0.0108; per speed ratio_of_minima 1.0); verification/paper/protocol_window.py/.json (summary: 4 points, 53 runs, n_disagreeing_runs 0, ratio 1.0); verification/paper/protocol_phases.py/.json (summary: n_phases_available 64; 1090 rpm min_16 10.881, min_32 10.763, min_64 10.763 mm/s, rel_change_of_min -0.0108, spread_16 0.139, spread_available 0.151; 1420, 2500, 2870 rpm rel_change 0, spread unchanged)

\subsection{Describing function and gain plane}\label{supp:bistab-df}

Equation~\eqref{M-eq:df} was evaluated as the average over one tooth period and over the phase of the chatter component, on $128\times128$ samples. The factor $p$ was fitted on the unclipped cut kicked by \SI{1}{\milli\metre\per\second} at the start of a tooth period: the velocity deviation from the unkicked run was high-passed at \SI{40}{\hertz}, which removes the actuator mode, and the per-tooth peak of the deviation, once the faster modes had died out, was fitted by the slowest mode and extrapolated back to the kick; fits with $r^2<0.98$ were discarded, which left $p$ from $0.7\alim$ upwards ($0.8\alim$ at \SI{1420}{\rpm}). On a machine, $p$ could be estimated from a measured frequency response as the residue of the critical root of Eq.~\eqref{M-eq:char}; this was not attempted.
% src: src/c1_unsafe_zone.py (df_gain: _TH 128 phases, b on ND = 128 samples per tooth period; _mode_factor: dv0 = 1e-3 m/s at the start of revolution 100, F_HP = 40 Hz, k0 = 6, r2 >= 0.98); results/c1_basin.npz (p_mode finite from 0.7 a_lim, 0.8 a_lim at 1420 rpm)

At \SI{1090}{\rpm} the clipped step-entry limit stays at \SIrange{19.3}{19.9}{\milli\metre} for $g\ge\SI{1650}{\volt\second\per\metre}$ (\SI{18.8}{\milli\metre} at $g=1500$), and at \SI{2500}{\rpm} at \SIrange{13.0}{13.2}{\milli\metre} for $g\ge1500$. The face entry follows $\alim$ until the forced command itself exceeds the clip at some depth below $\alim$: from $g=1650$ at \SI{2500}{\rpm}, after which its limit stays at \SIrange{16.7}{17.9}{\milli\metre}, and only at $g=3000$ at \SI{1090}{\rpm} (\SI{40.4}{\milli\metre} against $\alim=\SI{41.8}{\milli\metre}$). Below the plateau the clipped step-entry limit rises monotonically with the gain at both speeds; at \SI{2500}{\rpm} it is 6.4, 8.8 and \SI{13.0}{\milli\metre} at $g=536$, 800 and \SI{1500}{\volt\second\per\metre}.
% src: results/c1_gainplane.npz (gains 0-3000 in steps of 150; a_step_clip monotone in g at both speeds; a_edge_clip, a_lim, orbit_peak_V: 9.21 V at g = 1650, 2500 rpm; 9.16 V at g = 3000, 1090 rpm), results/c1_summary.json (gainplane); results/table3_main.csv (2500 rpm: a_prop_td 6.373 mm at g = 535.6, a_clip1500_td 13.010 mm), results/table3_clip_gain_sweep.csv (8.771 mm at g = 800); README §5.5, §11.6
